\begin{textsample}{POS Dim 3 – vem\_ed\_32 – Score 4.00 – vem\_ed\_32/t171.txt}  \label{ex:f3_pos_039}
Osvaldo Succi em webinário da UNICollaboration Osvaldo Succi, coordenador dos Projetos Colaborativos Internacionais (PCIs) na equipe de Apoio à Internacionalização do Ensino Superior da Divisão de Extensão e Pesquisa no Ensino Superior (Depes) da Coordenadoria Geral de Ensino Superior de Graduação (CGESG) do Centro Paula Souza, ministrou webinário a convite da UNICollaboration em 31 de \textbf{outubro} de 2025.
A organização europeia tem o objetivo de promover a integração entre pesquisa e prática de Intercâmbios Virtuais.
O título da apresentação foi “The use of Artificial Intelligence by multilingual speakers and VE project coordinators” (O uso da Inteligência Artificial por falantes multilíngues e coordenadores de projetos de Intercâmbios Virtuais).
Entre os principais pontos desfavoráveis ao uso da IA por falantes multilíngues, Osvaldo Succi \textbf{apontou} limitações de interação humana e de contexto cultural; dependência excessiva e autonomia reduzida; questões éticas e de privacidade; barreiras técnicas e inequidade; além de qualidade e consistência.
Por outro lado, a IA ajuda os coordenadores de Intercâmbios Virtuais a aumentar a criatividade ao parear professores de disciplinas diferentes; pular perguntas básicas e fazer questões mais inspiradoras sobre as áreas específicas de conhecimento dos docentes; considerar diferentes tipos de Intercâmbio Virtual (em especial os inter e transdisciplinares); manter melhores registros das reuniões; lidar com tarefas repetitivas, porém importantes, como detalhar instruções para os alunos e, assim, liberar tempo para os professores parceiros conversarem e fortalecerem seus laços de relacionamento.
Enquanto a inteligência humana é melhor para eticamente avaliar soluções, aplicando sabedoria e contexto, a artificial auxilia a estimular cenários e sugerir soluções \textbf{orientadas} por dados.
Se a IA proporciona acesso imediato a quantidades enormes de informação, por outro lado a qualidade humana da curiosidade, de desaprender e reaprender é única.
O importante é que sempre haja a mediação humana para somar e ponderar o melhor de cada uma das inteligências, com ética e \textbf{espírito} crítico.

% matched lemmas: apontar, espírito, orientar, outubro
\end{textsample}
